%=====================================================================
\chapter{Background and Preliminaries}
\label{ch:background}
%=====================================================================

This chapter assembles the technical material on which the rest of the thesis depends. It is
deliberately self-contained: a reader who works through it should be able to read
Chapters~\ref{ch:fusion} to~\ref{ch:allocation} without recourse to external references for
notation, for the definition of a metric or for the justification of a statistical procedure.
The chapter does not survey the research literature --- that is the task of
Chapter~\ref{ch:litsurvey} --- but establishes the objects the literature operates on.

\section{The Market Environment and Its Information Streams}

\subsection{Market efficiency and the scale of the achievable edge}

Any claim about forecastable structure in asset prices must be positioned against the
efficient market hypothesis in the form given by Fama~\cite{fama1970}: if prices reflect
available information, systematic outperformance net of costs should not persist. The
hypothesis is not falsified by the existence of a small measurable edge; it constrains its
size. The benchmark machine-learning study in empirical asset pricing, by Gu, Kelly and
Xiu~\cite{gu2020}, reports exactly this pattern --- statistically detectable predictability
whose magnitude is modest and whose economic value depends heavily on how it is harvested.

This framing is used throughout the thesis in a specific way. When Chapter~\ref{ch:llm}
reports a directional accuracy of 47.72 per cent against a baseline mean of 46.43 per cent,
the correct reaction is not disappointment but recognition that the measured increment is of
the order the literature leads one to expect at a single-name daily horizon. Conversely, when
Chapter~\ref{ch:fusion} reports 84.7 per cent at the index level, the correct reaction is to
ask what object is being predicted and over what horizon, because the two figures are not
comparable. Section~\ref{sec:reconcile} of Chapter~\ref{ch:conclusion} treats that
reconciliation as a substantive finding rather than as a caveat.

Two qualifications to the strict efficiency reading are relevant. The first is empirical: robust
cross-sectional regularities exist, of which the size and value factors of Fama and
French~\cite{famafrench1993} and the momentum effect of Jegadeesh and
Titman~\cite{jegadeesh1993} are the most durable, and their persistence establishes that
predictable structure is a matter of degree rather than of existence. The second is behavioural.
Prospect theory~\cite{kahneman1979} and the behavioural-finance programme that followed
it~\cite{shiller2003} describe systematic departures from the rationality that strict efficiency
assumes; Lo's adaptive markets hypothesis~\cite{lo2004} reconciles the two positions by treating
efficiency as an evolving property of a population of strategies rather than a fixed state. On
that reading, an edge can exist, be competed away and re-emerge --- which is precisely the
non-stationarity of source reliability that Chapter~\ref{ch:dynfusion} is built to accommodate.

\subsection{Structural characteristics of the Indian equity market}

The National Stock Exchange of India and the Bombay Stock Exchange together constitute the
primary venue for Indian equity trading. Three characteristics are relevant to the design
decisions taken in this thesis.

\textbf{Retail participation and sentiment transmission.} Retail participation in Indian
equities is high and has risen sharply over the study period. Sentiment expressed on social
platforms is therefore not only a commentary on the market but a partial determinant of
short-term price movement, and the relationship is asymmetric: Khan and
Chahal~\cite{khan2025asymmetric} find that the influence of social media sentiment on Indian
sectoral returns is strongest at market extremes. This motivates the inclusion of social media
as a distinct stream in Chapter~\ref{ch:fusion} rather than as an adjunct to news.

\textbf{Policy sensitivity and discrete structural breaks.} The annual Union Budget, monetary
policy committee announcements by the Reserve Bank of India, and regulatory decisions of the
Securities and Exchange Board of India produce discrete, scheduled events that reset
relationships between instruments within hours. This is the reason a \emph{static} causal
graph is inadequate for the Indian market and a time-varying one is required, as
Chapter~\ref{ch:causal} argues; it is also the reason the stationarity assumption embedded in
standard debiasing layers fails at exactly the moments that matter.

\textbf{Sectoral concentration and dual exposure.} The index carries pronounced concentration
in financial services, so shocks in that sector propagate to the aggregate disproportionately;
and the market is simultaneously exposed to global factors, including crude oil prices, the
rupee--dollar exchange rate and United States monetary policy. Domestic and international
information can therefore disagree, and the disagreement is informative. This is what makes
the market a demanding test of a fusion architecture rather than a convenient one.

\subsection{The four information streams}

Table~\ref{tab:streams} summarises the four streams used in this thesis, their native sampling
frequency, the transformation applied to reach a common daily representation, and the
distortion that transformation introduces. The final column is the important one: every
alignment step buys comparability at the cost of some fidelity, and stating the cost is part of
the methodology rather than an afterthought.

\begin{table}[!htb]
\centering
\caption[The four information streams and their alignment costs]{The four information streams,
their native frequency, the transformation applied to reach a common daily representation, and
the distortion that transformation introduces}
\label{tab:streams}
\tabfont
\begin{tabular}{@{}L{2.3cm}L{2.4cm}L{4.3cm}L{4.3cm}@{}}
\toprule
\textbf{Stream} & \textbf{Native frequency} & \textbf{Transformation to daily} &
\textbf{Distortion introduced} \\
\midrule
Market data (OHLCV and derived indicators) & Per session, and intraday where available &
Indicators computed on closing series; intraday series aggregated to session features & None
at the daily horizon; intraday aggregation discards microstructure ordering \\
News text & Irregular, time-stamped & Scored per article, aggregated per session, with
post-close publications attributed to the next session & Within-day ordering of articles is
lost; relevance weighting is a modelling choice, not a measurement \\
Social media text & Irregular, high volume & Scored per post, aggregated into daily polarity,
mention volume and dispersion & Bot and coordinated-posting contamination is not separated;
dispersion partly compensates \\
Macroeconomic series & Monthly or quarterly & Interpolated to daily frequency &
Announcement-day effects are discontinuous and interpolation understates them; the direction
of the resulting bias is known and stated \\
\bottomrule
\end{tabular}
\end{table}

\section{Technical Indicators}

Technical indicators compress the price and volume record into features intended to expose
trend, momentum, volatility and participation. Fifteen indicators are used in
Chapters~\ref{ch:fusion} to~\ref{ch:llm}; the families and the reasons for their inclusion are
given here, with the canonical definitions, which follow the standard
references~\cite{murphy1999,wilder1978,bollinger2001}. No claim is made that any individual
indicator carries independent predictive power; they are used as a feature basis that the
learning architecture combines, and the question of which of them matters is answered by the
ablations of Chapter~\ref{ch:fusion} rather than assumed here.

\subsection{Trend and moving averages}

The simple moving average over $n$ sessions,
\begin{equation}
\mathrm{SMA}_n(t) = \frac{1}{n}\sum_{k=0}^{n-1} P_{t-k},
\end{equation}
is the elementary trend statistic; the exponential moving average,
$\mathrm{EMA}_n(t) = \beta P_t + (1-\beta)\,\mathrm{EMA}_n(t-1)$ with
$\beta = 2/(n+1)$, weights recent observations more heavily and responds faster to a turn.
Windows of 5, 10 and 20 sessions are used, capturing weekly, fortnightly and monthly trend
respectively. The slope of the exponential average supplies a directional feature that the
level does not.

\subsection{Momentum oscillators}

The relative strength index of Wilder~\cite{wilder1978} normalises the ratio of average gains
to average losses over $n$ sessions onto a bounded scale,
\begin{equation}
\mathrm{RSI}_n(t) = 100 - \frac{100}{1 + \mathrm{RS}_n(t)},
\qquad
\mathrm{RS}_n(t) = \frac{\overline{\text{gain}}_n(t)}{\overline{\text{loss}}_n(t)},
\end{equation}
and is conventionally read as indicating overbought conditions above 70 and oversold conditions
below 30. The moving average convergence divergence statistic is the difference between two
exponential averages, $\mathrm{MACD}(t) = \mathrm{EMA}_{12}(t) - \mathrm{EMA}_{26}(t)$, with a
signal line given by $\mathrm{EMA}_{9}$ of the statistic itself; the crossing of the two is the
conventional trigger. The Aroon indicators measure the number of sessions since the highest
high and the lowest low within a window, and the average directional index summarises trend
strength independently of direction.

\subsection{Volatility and range}

Bollinger bands place a $k$-standard-deviation envelope around a moving average,
$\mathrm{SMA}_n(t) \pm k\,\sigma_n(t)$, with $k=2$ conventional; the band width and the
position of the close within the band are both used as features. The average true range
averages the maximum of the session range and the gaps to the previous close, and realised
volatility over a window is computed as the standard deviation of log returns
$r_t = \log(P_t / P_{t-1})$, annualised by $\sqrt{252}$.

That volatility is itself predictable, clusters in time, and responds asymmetrically to negative
returns is the standard description of financial return series since the autoregressive
conditional heteroskedasticity model of Engle~\cite{engle1982} and its generalisation by
Bollerslev~\cite{bollerslev1986} --- a description that the linear
autoregressive-integrated-moving-average machinery of Box and Jenkins~\cite{boxjenkins1970},
which supplies the classical baseline in Chapter~\ref{ch:fusion}, cannot represent. This matters
twice in the thesis: it is the reason a regime-conditional risk budget is worth constructing at
all (Chapter~\ref{ch:allocation}), and it is the reason a debiasing layer that assumes stationary
noise fails under stress (Chapter~\ref{ch:causal}).

\subsection{Participation and volume}

On-balance volume accumulates signed volume,
$\mathrm{OBV}(t) = \mathrm{OBV}(t-1) + \mathrm{sgn}(P_t - P_{t-1})\,V_t$, and the
volume-weighted average price,
$\mathrm{VWAP}(t) = \sum_k P_k V_k / \sum_k V_k$, supplies a participation-weighted reference
level. Volume momentum, the ratio of current volume to its own moving average, distinguishes a
price move made on conviction from one made on thin trading.

A caution applies to the unbounded accumulators. On-balance volume grows without bound and
takes values of the order $10^{8}$ to $10^{9}$ on liquid instruments, while a risk term in a
mean-variance objective is of the order $10^{-4}$. Passing such a feature into an optimiser
without rescaling lets the ratio of units, rather than the information in the signal, determine
the solution. This is not a hypothetical concern: it is one of the two engineering failure
modes diagnosed in Chapter~\ref{ch:allocation}.

\section{Sentiment Extraction from Financial Text}

Financial text analysis has passed through three generations, and the distinctions between them
determine what a sentiment score means.

\subsection{Lexicon methods}

The first generation scores text against a predefined polarity dictionary. The decisive finding
for finance is that of Loughran and McDonald~\cite{loughran2011}: general-purpose lexicons
misclassify financial language badly, because words that are negative in ordinary usage
(``liability'', ``tax'', ``cost'') are neutral descriptors in a filing. Domain-specific
dictionaries are therefore required. VADER~\cite{hutto2014} remains a widely used reference
implementation for social text, augmenting a lexicon with rules for intensifiers, negation,
capitalisation and punctuation. Lexicon methods are fast, fully interpretable and deterministic;
they are insensitive to context, sarcasm and syntax.

\subsection{Supervised machine learning}

The second generation learns polarity from annotated corpora using conventional classifiers.
Accuracy improves over lexicon baselines but at the cost of annotation, and generalisation
across market conditions is weak: a classifier trained on one regime's vocabulary degrades when
the vocabulary shifts~\cite{rouf2021}.

\subsection{Transformer language models}

The third generation replaced both. The transformer architecture~\cite{vaswani2017} computes
representations through scaled dot-product self-attention,
\begin{equation}
\mathrm{Attention}(Q,K,V) = \mathrm{softmax}\!\left(\frac{QK^{\!\top}}{\sqrt{d_k}}\right) V,
\end{equation}
where $Q$, $K$ and $V$ are query, key and value projections of the input and $d_k$ is the key
dimension. The same mechanism, applied across modalities rather than across tokens, is the
cross-modal attention used for fusion in Chapter~\ref{ch:fusion}; the notational continuity is
deliberate.

The bidirectional pre-training objective of BERT~\cite{devlin2019bert} and the optimised
training recipe of RoBERTa~\cite{liu2019roberta} established the encoder family on which
financial adaptation is built, while the demonstration that sufficiently large autoregressive
models perform new tasks from instructions alone~\cite{brown2020gpt3} opened the generative route
to sentiment extraction that Chapter~\ref{ch:llm} does not take but that
Section~\ref{sec:llm-rationale} discusses.

Domain adaptation followed quickly. FinBERT~\cite{araci2019} further-trains a bidirectional
encoder on financial text; larger domain models~\cite{wu2023bloomberggpt,yang2023fingpt,
wang2024fingptbench} extend the approach to generative reasoning over financial documents, and
surveys of their practical application record both the gains and the operational
obstacles~\cite{khan2024chatgpt}. Jung and Jang~\cite{jung2024} showed that masking numerically
significant tokens during adaptation preserves quantitative context that generic pre-processing
destroys --- ``GDP growth at 7.8 per cent'' carries information that ``GDP growth at [NUM] per
cent'' does not.

\subsection{Continuous versus categorical scores}

A design decision recurs throughout this thesis: sentiment is retained as a continuous polarity
score $s_t \in [-1,+1]$ rather than as a categorical label. Categorical output discards
intensity, and intensity is exactly what distinguishes a routine positive report from a
significant one. Retaining the continuous score, together with the model's own confidence,
preserves that information for the fusion layer and makes it possible to examine the empirical
distribution of sentiment as a diagnostic --- a symmetric distribution about neutral with
moderate tails indicates that the pipeline is measuring variation rather than producing a biased
categorical output, and this check is reported in Chapter~\ref{ch:llm}.

\section{Multi-Source Data Fusion}

\subsection{The three canonical strategies}

Given $M$ heterogeneous sources with per-source representations $x_1, \dots, x_M$ and a target
$y$, three fusion strategies dominate.

\textbf{Early (feature-level) fusion} concatenates representations before modelling,
$\hat{y} = f([x_1; \dots; x_M])$. It is simple and preserves cross-feature interactions, but
forces one architecture to learn interactions across domains with very different statistical
properties, and is prone to modality dominance if scales are not equalised~\cite{moodi2024}.

\textbf{Late (decision-level) fusion} trains source-specific models and combines their outputs,
$\hat{y} = \sum_m w_m f_m(x_m)$. It preserves modularity and allows each architecture to match
its source, but the weights $w_m$ are typically fixed in advance or learned once and then
frozen~\cite{farimani2024,vargas2022}.

\textbf{Hybrid (model-level) fusion} combines intermediate representations inside the network,
usually through attention or gating~\cite{liu2019attention,haryono2023,zong2025}. This is the
current state of the art and is a genuine advance, because the combination is learned rather
than assumed.

\subsection{What attention weights do and do not express}

The limitation that motivates Chapter~\ref{ch:dynfusion} can be stated precisely. A learned
attention weight $\alpha_m$ is a point estimate obtained by minimising a loss over an entire
training set. It expresses the \emph{average} relevance of modality $m$ across that set. It is
not a statement about the reliability of modality $m$ at a particular moment, and it carries no
uncertainty. Consequently, if a source degrades after training --- because a news feed changes
composition, because a social channel is contaminated, or because a regime shift alters which
signals carry information --- the architecture continues to weight it at its historical value
until the next retraining cycle.

The remedy pursued in this thesis is to make the weight a function of a \emph{measured}
quantity. If $e_{i,t} = y_t - \hat{y}_{i,t}$ is the prediction error of source $i$, then
\begin{equation}
\sigma_{i,t}^2 = \frac{1}{\tau}\sum_{k=1}^{\tau}\left(e_{i,t-k} - \bar{e}_{i,t}\right)^2
\label{eq:uncertainty}
\end{equation}
is an empirical estimate of that source's recent predictive volatility over a trailing window
of $\tau$ sessions, and the fusion weights
\begin{equation}
w_{i,t} = \frac{\exp\!\left(-\sigma_{i,t}^2 / T\right)}
               {\sum_j \exp\!\left(-\sigma_{j,t}^2 / T\right)}
\label{eq:softmaxweights}
\end{equation}
discount a source in proportion to how unreliable it has recently been, with temperature $T$
controlling the sharpness of the discounting. The estimate in Equation~\eqref{eq:uncertainty} is
\emph{aleatoric}: it captures variation attributable to noise in the environment. It does not
capture \emph{epistemic} uncertainty --- the model's ignorance about regions of input space it
has not seen --- for which a Bayesian treatment of the network weights would be
required~\cite{gal2016dropout,kendall2017}. The distinction is maintained in
Chapter~\ref{ch:dynfusion} and the epistemic extension is identified as future work.

\subsection{Distribution-aware normalisation}

Heterogeneous sources have heterogeneous distributions, and a single scaler applied to all of
them lets the heaviest-tailed input dominate the fused representation. Three scalers are used
according to the shape of the input: standard scaling, $(x - \mu)/\sigma$, for approximately
symmetric series such as log returns; min--max scaling, $(x - x_{\min})/(x_{\max} - x_{\min})$,
for bounded series such as the sentiment score, which preserves the interpretable range; and
robust scaling, $(x - \mathrm{med})/\mathrm{IQR}$, for heavy-tailed series such as a volatility
index, where the mean and standard deviation are themselves unstable.

Every scaler in this thesis is fitted on an expanding window of strictly earlier observations.
Fitting a scaler on the full sample is a subtle and common form of look-ahead: the test-period
mean and variance leak into the training representation even though no test label is used.

\section{Temporal Alignment and Look-Ahead}

The streams of Table~\ref{tab:streams} arrive on different clocks, and reconciling them is
where look-ahead bias most often enters. Four rules are applied throughout this thesis.

\textbf{Publication-time attribution.} A news item or post is attributed to the first trading
session that begins strictly after its time stamp. An article published after the close is
attributed to the next session, never to the session that has already ended.

\textbf{Holidays are gaps, not repeats.} Market holidays are treated as genuine absences rather
than forward-filled. Forward-filling a closed session injects artificial autocorrelation into
the return series and inflates any statistic that depends on persistence.

\textbf{Macroeconomic release lag.} A macroeconomic series is available only after its release
date, not after the period it describes. Interpolation to daily frequency is applied between
releases, with the consequence --- understatement of announcement-day discontinuities ---
stated explicitly.

\textbf{Strictly chronological partitioning.} Training, validation and test blocks are
contiguous and ordered in time. Random or stratified splitting of a time series is a defect,
not a variant, because it permits the model to learn from the future of the observation it is
scored on.

\section{Learning Architectures}

\subsection{Recurrent networks}

Long short-term memory~\cite{hochreiter1997} maintains a cell state regulated by input, forget
and output gates, which permits gradients to propagate across long sequences without vanishing
--- the failure mode that Bengio, Simard and Frasconi~\cite{bengio1994} showed to be intrinsic to
simple recurrent networks trained by gradient descent. The bidirectional variant runs two such
networks in opposite directions and concatenates their hidden states,
$h_t = [\overrightarrow{h}_t; \overleftarrow{h}_t]$, so that the representation at each step
reflects both past and future context \emph{within the input window}. This is not look-ahead
provided the input window closes at or before the decision time; the distinction is between
context inside a supplied window and information from beyond it. A temporal attention layer over
the hidden sequence, in the form introduced for sequence alignment by Bahdanau, Cho and
Bengio~\cite{bahdanau2015}, allows the output to weight positions in that window unequally.

The gated recurrent unit~\cite{cho2014gru} merges the forget and input gates into a single
update gate and dispenses with the separate cell state. It has fewer parameters than long
short-term memory and trains faster on comparable data, which is why it is the market-series
expert in Chapter~\ref{ch:dynfusion} while the full bidirectional variant is used for the fused
representation in Chapter~\ref{ch:fusion}.

\subsection{Convolutional encoders for indicator streams}

One-dimensional convolution applied along the time axis of an indicator stream extracts local
patterns --- a crossover, a squeeze, a volume spike --- with far fewer parameters than a dense
layer, and with translation invariance that is appropriate because the same pattern carries the
same meaning wherever it occurs. The encoder used in Chapter~\ref{ch:fusion} is
\begin{equation}
h_m^{(l)} = \mathrm{ReLU}\!\left(\mathrm{Conv1D}\!\left(W_m^{(l)} * h_m^{(l-1)} +
b_m^{(l)}\right)\right),
\label{eq:conv}
\end{equation}
with $h_m^{(0)}$ the raw feature matrix of modality $m$, and max-pooling between layers.

\subsection{Gradient-boosted trees}

Gradient boosting, in the general formulation of Friedman~\cite{friedman2001}, fits an additive
ensemble of regression trees, each fitted to the gradient of the loss with respect to the current
prediction. The implementation used here~\cite{chen2016xgboost} adds second-order optimisation and
explicit regularisation on tree complexity; histogram-based alternatives~\cite{ke2017lightgbm}
trade a small amount of fidelity for speed and are not required at the sample sizes involved.
Tree ensembles are used in Chapter~\ref{ch:llm} rather than a deep network for three reasons:
they are invariant to monotone rescaling of features, which matters when combining indicators
measured in incompatible units; they handle heterogeneous feature types without embedding layers;
and their split-gain statistics supply a feature-importance measure that is intrinsic to the
fitted model rather than a post-hoc approximation. Random forests~\cite{breiman2001} and support
vector machines~\cite{cortes1995svm}, in the standard implementations of
Pedregosa et al.~\cite{pedregosa2011}, serve as baselines.

\subsection{Regularisation and optimisation}

Dropout~\cite{srivastava2014dropout} randomly zeroes activations during training, which prevents
co-adaptation and acts as an implicit ensemble; it is also the mechanism by which an approximate
predictive distribution can be obtained at inference time, as Section~2.4.2 noted. Batch
normalisation~\cite{ioffe2015bn} stabilises the distribution of layer inputs and permits larger
learning rates. Adaptive moment estimation~\cite{kingma2015adam} is used as the optimiser
throughout, with early stopping on validation loss and a fixed patience. Random seeds are fixed
for data shuffling and weight initialisation in every reported experiment, so that a difference
between two configurations reflects the architecture rather than the initialisation.

\section{Regime Modelling}

\subsection{Why regimes}

Financial series are not homogeneous through time. The insight formalised by
Hamilton~\cite{hamilton1989} is that a series may be generated by a small number of distinct
states with different means, variances and persistence, and that identifying the current state
is a precondition for acting sensibly. Ang and Bekaert~\cite{ang2002} demonstrated that
accounting for regime shifts materially changes optimal allocations, and Kritzman et
al.~\cite{kritzman2012} made the practical case for regime-aware dynamic strategies.

\subsection{Gaussian mixture identification}

The regime model used in Chapter~\ref{ch:allocation} is a $K$-component Gaussian mixture over a
two-dimensional feature vector of rolling realised volatility and rolling momentum. The mixture
density is
\begin{equation}
p(x) = \sum_{k=1}^{K} \pi_k \,\mathcal{N}\!\left(x \mid \mu_k, \Sigma_k\right),
\qquad \sum_k \pi_k = 1,
\end{equation}
fitted by expectation maximisation~\cite{dempster1977}, and the state assigned to a session is
the component with the highest posterior responsibility. Three components are used, interpreted
after anchoring as calm, transitional and stress states; the number is fixed \emph{a priori} on
interpretability grounds rather than selected by an information criterion such as
Schwarz's~\cite{schwarz1978}, because selecting $K$ on the evaluation sample would be a further
specification decision requiring disclosure, and three states map onto a distinction a
practitioner already makes.

Two implementation points are consequential and are stated here because they generalise.

\textbf{Label anchoring is mandatory.} The component indices returned by expectation
maximisation are arbitrary and permute freely between refits. If the mixture is refitted at
every rebalance --- which it must be, to avoid look-ahead --- then component~0 in one refit need
not correspond to component~0 in the next. Components must therefore be sorted by a stable,
interpretable statistic; ascending fitted volatility is used here. Omitting this step produces
no error message, only a silently scrambled regime series, and the downstream results will look
plausible while being meaningless.

\textbf{Refitting must use only earlier data.} The mixture is refitted at each rebalance on an
expanding window of strictly earlier sessions. Fitting once on the full sample and then
labelling the history is a severe and easily overlooked form of look-ahead, because the
identification of ``stress'' is informed by knowing when stress occurred.

\section{Causal Discovery}

\subsection{Predictive precedence}

The oldest operational notion of directed influence in time series is that of
Granger~\cite{granger1969}: series $X$ is said to Granger-cause series $Y$ if past values of
$X$ improve the prediction of $Y$ beyond what $Y$'s own past achieves. Formally, comparing the
residual variances of the restricted and unrestricted regressions
\begin{equation}
Y_t = \sum_{k=1}^{p}a_k Y_{t-k} + \varepsilon_t
\qquad\text{versus}\qquad
Y_t = \sum_{k=1}^{p}a_k Y_{t-k} + \sum_{k=1}^{p}b_k X_{t-k} + \eta_t,
\end{equation}
the null $b_1 = \dots = b_p = 0$ is tested by an $F$ statistic. The notion is predictive rather
than interventional: it establishes precedence in an information set, not a manipulable cause,
and the distinction between the two is the subject of the structural-causal
literature~\cite{pearl2009,peters2017}. The thesis uses the term ``causal'' in this operational
sense throughout and does not claim interventional identification. Two practical cautions from
the time-series literature apply and are observed: inference in a vector autoregression with
possibly integrated series requires the lag augmentation of Toda and
Yamamoto~\cite{toda1995} if the usual asymptotics are to hold, and the identification problems
that Sims~\cite{sims1980} raised for structural interpretation of such systems are not solved by
the procedure used here. Modern constraint-based methods designed for large nonlinear time-series
datasets~\cite{runge2019} address the scalability limitation discussed below.

\subsection{Constraint-based structure learning}

Constraint-based discovery procedures of the kind formalised by Spirtes, Glymour and
Scheines~\cite{spirtes2000} recover a directed structure by testing conditional independences.
Starting from a complete undirected graph, an edge between $X_i$ and $X_j$ is removed whenever a
conditioning set $S$ is found such that $X_i \perp\!\!\!\perp X_j \mid S$; the surviving skeleton
is then oriented using the conditioning sets that produced the removals. Applied to lagged
financial series, the procedure yields a directed adjacency matrix $A(t)$ whose entry $(i,j)$
records the strength of directed influence of series $i$ on series $j$.

The computational cost is the binding practical constraint: the number of conditional
independence tests grows combinatorially with the number of variables. This is why the study in
Chapter~\ref{ch:causal} uses five instruments rather than fifty, and why scalable discovery is
identified as a prerequisite for production-scale causal transparency.

\subsection{Correlation is not causality, and the difference is exploitable}

Two series can be strongly linearly correlated with neither influencing the other, and --- more
usefully for this thesis --- one series can exert strong directed influence on another while
their linear correlation is weak, if the influence operates with a lag or through a nonlinear
channel. A correlation-based feature selector discards exactly those relationships. The
empirical demonstration that such pairs exist in the studied universe is the mechanistic
justification for the causal approach, and it is reported in Chapter~\ref{ch:causal}.

\section{Portfolio Construction}

\subsection{Mean-variance and its estimation problem}

The mean-variance formulation of Markowitz~\cite{markowitz1952} selects weights $w$ maximising
$\mu^{\!\top} w - \lambda\, w^{\!\top}\Sigma w$ for expected returns $\mu$, covariance $\Sigma$
and risk aversion $\lambda$. The formulation is correct and the estimation is fragile: sample
estimates of $\mu$ and $\Sigma$ carry large errors that the optimiser amplifies, because it
allocates aggressively to whatever the estimates most favour --- the ``optimization enigma''
described by Michaud~\cite{michaud1989}, who characterised mean-variance optimisers as
error-maximisers. Four responses matter here. Ledoit and Wolf~\cite{ledoit2004} shrink the sample
covariance toward a structured target, $\hat{\Sigma} = (1-\delta)S + \delta F$, producing a
well-conditioned estimator at the cost of a small bias. Black and
Litterman~\cite{blacklitterman1992} anchor expected returns on equilibrium and blend views into
them, which stabilises $\mu$ rather than $\Sigma$. Jagannathan and Ma~\cite{jagannathan2003}
showed that a long-only constraint of the kind used in Chapter~\ref{ch:allocation} is itself a
form of shrinkage and improves out-of-sample performance for that reason. And risk-parity
constructions~\cite{maillard2010} dispense with expected returns altogether, equalising risk
contributions. DeMiguel, Garlappi and Uppal~\cite{demiguel2009} demonstrated that naive equal
weighting frequently outperforms optimisation that ignores estimation error --- which is why
equal weighting is used as the primary benchmark in Chapter~\ref{ch:allocation} rather than as a
straw man. Ang~\cite{ang2014} provides the systematic treatment of the factor-based alternative.

\subsection{Converting a signal into an expected return}

A predictive signal is not an expected return and cannot be substituted for one. The standard
conversion is information-ratio scaling in the form given by Grinold and
Kahn~\cite{grinold2000}: a cross-sectionally standardised score $z_{i,t}$ becomes an expected
excess return through
\begin{equation}
\alpha_{i,t} = \mathrm{IC} \cdot \sigma_{i,t} \cdot z_{i,t},
\label{eq:grinold}
\end{equation}
where $\mathrm{IC}$ is the assumed information coefficient of the signal and $\sigma_{i,t}$ the
asset's forecast volatility. The scaling matters for a reason that is not cosmetic: the linear
term $\alpha^{\!\top}w$ and the quadratic term $w^{\!\top}\Sigma w$ in the objective must be in
compatible units, or their ratio --- rather than the information in the signal --- determines the
solution. Chapter~\ref{ch:allocation} reports what happens when this is neglected.

\subsection{Position sizing under an estimated edge}

The growth-optimal criterion of Kelly~\cite{kelly1956} gives the bet size maximising long-run
log wealth; for continuous returns with signal strength $s$ and forecast variance
$\sigma^2$, $f^{*} = \mathrm{IC}\cdot s / \sigma^2$. Full Kelly is rarely appropriate in
practice, because the criterion assumes $\mathrm{IC}$ and $\sigma$ are known rather than
estimated; a fractional multiple is standard, and the prototype system described in
Chapter~\ref{ch:allocation} uses one quarter of full Kelly with per-name and per-sector caps.

\subsection{Risk budgeting: relative rather than absolute}

Volatility-managed portfolios~\cite{moreira2017} scale total exposure by recent volatility.
Chapter~\ref{ch:allocation} adopts the idea but with one modification that proves decisive: the
risk budget is expressed as a \emph{fraction of the portfolio's own realised volatility} rather
than as an absolute annualised target. An absolute target of, say, 15 per cent is not scale-free;
it is calibrated for a concentrated equity sleeve and will essentially never bind on a
diversified multi-asset sleeve whose own volatility is 6 per cent. The layer is then present in
the specification, consumes computation, and does nothing. A relative budget is scale-free and
remains active in any universe.

\section{Evaluation Metrics}

\subsection{Forecast quality}

For predicted values $\hat{y}_i$ and realised values $y_i$ over $n$ observations,
\begin{equation}
\mathrm{MAE} = \frac{1}{n}\sum_{i=1}^{n}\left|y_i - \hat{y}_i\right|,
\qquad
\mathrm{RMSE} = \sqrt{\frac{1}{n}\sum_{i=1}^{n}\left(y_i - \hat{y}_i\right)^2},
\end{equation}
\begin{equation}
R^2 = 1 - \frac{\sum_{i}(y_i - \hat{y}_i)^2}{\sum_{i}(y_i - \bar{y})^2},
\qquad
\mathrm{DA} = \frac{1}{n}\sum_{i=1}^{n}
\mathbf{1}\!\left[\mathrm{sgn}(\Delta y_i) = \mathrm{sgn}(\Delta \hat{y}_i)\right].
\end{equation}
Where errors are reported as normalised, prices have been min--max scaled before the
computation, so the reported magnitudes are dimensionless and comparable across instruments.
Directional accuracy must always be read against the base rate of the majority class; on
approximately balanced daily equity direction that base rate is close to one half, and this is
the correct reference in Chapter~\ref{ch:llm}.

\subsection{Classification quality}

Precision, recall and the $F_1$ score are reported alongside accuracy in
Chapter~\ref{ch:llm} because accuracy alone conceals degenerate behaviour: a classifier that
predicts the positive class unconditionally attains perfect recall and uninformative precision,
and the reported results contain exactly such a case. Receiver operating characteristic analysis
summarises the trade-off across thresholds; an area under the curve close to 0.5 indicates
performance indistinguishable from random ranking.

\subsection{Portfolio quality}

For a return series $r_t$ with risk-free rate $r_f$,
\begin{equation}
\mathrm{Sharpe} = \frac{\overline{r} - r_f}{\sigma_r}\sqrt{252},
\end{equation}
following Sharpe~\cite{sharpe1966}. Maximum drawdown is the largest peak-to-trough decline of the
cumulative wealth series, $\mathrm{MDD} = \min_t ( W_t / \max_{s \le t} W_s - 1 )$, and the Calmar
ratio is the compound annual growth rate divided by the absolute maximum drawdown. Downside and
upside capture measure the fraction of benchmark movement participated in when the benchmark
falls and rises respectively; their ratio distinguishes genuine asymmetry from uniform
de-risking, and Chapter~\ref{ch:allocation} uses exactly this distinction.

Because the Sharpe ratio penalises upside and downside deviation alike, downside-sensitive
alternatives exist --- the Sortino ratio~\cite{sortino1994} divides excess return by downside
deviation only --- and the tail-risk measure used in the credibilistic optimiser of
Chapter~\ref{ch:causal} is conditional value at risk in the form of Rockafellar and
Uryasev~\cite{rockafellar2000}, which unlike value at risk satisfies the coherence axioms of
Artzner et al.~\cite{artzner1999} and is therefore well behaved under aggregation. The primary
reporting in this thesis nonetheless uses Sharpe ratio and maximum drawdown, because they are
what the comparison literature reports and substituting a more favourable measure would obscure
rather than clarify.

\subsection{Ablation at matched exposure}

A methodological point is recorded here because it is used as an instrument rather than merely
cited. Comparing a de-risked strategy against a fully invested benchmark confounds two effects:
the strategy holds \emph{less} risk on average, and it holds risk at \emph{different times}. The
first effect alone will improve a drawdown statistic and may improve a Sharpe ratio, without any
skill in timing. The remedy is a matched-exposure control: a variant carrying the same average
invested exposure as the full framework but allocating it uniformly rather than by regime. Any
difference between the framework and that control is attributable to the timing of risk alone.
This construction is possible only because the framework separates relative composition from
total risk by design, and it is the central evaluation instrument of
Chapter~\ref{ch:allocation}.

\section{Statistical Procedures}

\subsection{Inference under serial dependence}

Daily financial returns are serially dependent and heteroskedastic, so ordinary standard errors
understate sampling variability. The Newey--West estimator~\cite{newey1987} produces a
heteroskedasticity- and autocorrelation-consistent covariance matrix and is used for every
$t$ statistic on a time series in this thesis. For statistics whose sampling distribution is not
analytically tractable --- a Sharpe ratio difference, for example --- the stationary
bootstrap~\cite{politis1994} resamples blocks of random geometric length, preserving short-range
dependence, and the reported interval is the empirical quantile range across resamples.

\subsection{Comparing forecasts}

The Diebold--Mariano test~\cite{diebold1995} compares the predictive accuracy of two forecasts
by testing whether the mean of the loss differential $d_t = L(e_{1,t}) - L(e_{2,t})$ is zero,
using a long-run variance estimate for the standard error. It is used in Chapter~\ref{ch:llm} to
establish whether an observed accuracy difference between the sentiment-augmented model and its
baselines is distinguishable from sampling noise.

\subsection{Multiple testing and the credibility of a reported gain}

The most important statistical point in this thesis is not a test but a disclosure discipline.
Harvey and Liu~\cite{harvey2015} showed that conventional significance thresholds are inadequate
once the number of strategies examined is taken into account, and Bailey and López de
Prado~\cite{bailey2014dsr} formalised the correction through the deflated Sharpe ratio, which
adjusts for the number of trials, the non-normality of returns and the length of the sample.
Formal multiple-comparison procedures for exactly this setting exist --- White's reality
check~\cite{white2000} and Hansen's test for superior predictive ability~\cite{hansen2005} both
construct a null distribution over the full set of candidate strategies rather than over the
selected one. The practical consequence adopted here is simpler and stricter: the \emph{entire}
specification sweep is reported rather than its best cell. Chapter~\ref{ch:allocation}
accordingly reports all eighteen cells of a three-seed by two-covariance-window by
three-rebalance-interval grid, and states the range rather than the maximum.

\subsection{Robustness and significance are different properties}

A final distinction, applied consistently: a result is \emph{robust} if it holds across
specifications, sub-periods and universes; it is \emph{significant} if the observed effect is
unlikely under a null of no effect. These are different properties and neither implies the
other. Eighteen years of daily data contain thousands of observations but only a handful of
independent crisis episodes, so a drawdown claim can be highly robust and simultaneously fail a
significance test, because the effective sample size for the events that matter is small.
Chapter~\ref{ch:allocation} reports both and does not conflate them.

\section{Summary}

This chapter has established the objects the thesis operates on: the four information streams
and the cost of aligning them; the indicator families and the scale hazard that unbounded
accumulators introduce; the three generations of sentiment extraction and the reason a
continuous score is retained; the three fusion strategies and the precise limitation of learned
attention that motivates uncertainty-weighted fusion; the learning architectures and the
regularisation protocol; regime identification and the two implementation points that make it
valid; causal discovery in the predictive-precedence sense and the reason correlational
selection destroys usable information; portfolio construction, signal scaling and the case for a
relative risk budget; the evaluation metrics, including the matched-exposure ablation that is
this thesis's principal analytical instrument; and the statistical procedures, including the
disclosure discipline that makes a reported gain credible. Chapter~\ref{ch:litsurvey} now turns
to what the research literature has established using these objects, and to what it has left
unresolved.
